\section{Engine Solver System}
\label{sec:engine-solver-system}

This section documents the solver infrastructure under
\texttt{include/hacdcpf/engine} and \texttt{src/engine} as a coherent system.
The goal is to separate three concerns that are easy to blur when reading the
repository file-by-file:
\begin{itemize}
  \item the \emph{canonical mathematical problem forms} that every backend sees,
  \item the \emph{adapter and registry layer} that dispatches a problem to a solver,
  \item the \emph{actual kernels} that are fully implemented today versus declared extension points.
\end{itemize}

This chapter is intentionally code-aligned.  When a header advertises a solver
family but the current source file is only a placeholder, the notebook states
that explicitly instead of describing the ideal future design.

\subsection{Layering Inside \texttt{hacdcpf::engine}}

The engine layer is organized in five contracts.
\begin{enumerate}
  \item \textbf{Canonical problem layer:} \texttt{problem\_types.hpp} defines the shared mathematical models such as \texttt{LPModel}, \texttt{QPModel}, and \texttt{MIPModel}.
  \item \textbf{Validation layer:} \texttt{problem\_validation.hpp/.cpp} checks structural consistency before the numerical solve begins.
  \item \textbf{Adapter layer:} \texttt{solver\_adapter.hpp} defines the uniform interface \texttt{solve\_lp}, \texttt{solve\_nlp}, \texttt{solve\_milp}, and so on.
  \item \textbf{Registry layer:} \texttt{adapter\_registry.hpp/.cpp} stores candidate adapters and selects one by supported \texttt{ProblemClass} or by name.
  \item \textbf{Kernel layer:} concrete native and external implementations live in files such as \texttt{dual\_simplex.cpp}, \texttt{lcqp\_solver.cpp}, \texttt{pdlp\_solver.cpp}, and \texttt{external\_adapters.cpp}.
\end{enumerate}

Two different solve paths coexist inside this layer:
\begin{itemize}
  \item \textbf{Canonical optimization path:} a caller builds a \texttt{ProblemClass} object and dispatches it through a \texttt{SolverAdapter}.
  \item \textbf{Direct application kernel path:} some modules call engine kernels directly without going through the adapter abstraction.  The most important example is \texttt{NewtonSolver}, which consumes \texttt{powerflow::SolverData} rather than \texttt{NonlinearSystem}.
\end{itemize}

\subsection{Canonical Problem Classes and Standard Forms}

Table~\ref{tab:engine-problem-classes} is the central contract of the engine.
Every solver family either consumes one of these objects directly or wraps a
domain-specific model into one of these forms first.

\begin{longtable}{L{0.10\linewidth}L{0.17\linewidth}L{0.37\linewidth}L{0.26\linewidth}}
\toprule
Class & C++ struct & Standard mathematical form & Current engine role \\
\midrule
\endhead
\texttt{LE} & \texttt{SparseLinSys} & Solve
\(A x = b\). & Small reusable sparse linear-system abstraction used by \texttt{NativeLinearAdapter}. \\
\texttt{NLE} & \texttt{NonlinearSystem} & Find
\(F(x)=0\), with residual and Jacobian callbacks and optional projection. & Generic Newton-style nonlinear equation interface. \\
\texttt{LP} & \texttt{LPModel} &
\(\begin{aligned}
\min/\max \quad & c^\top x \\
\text{s.t.}\quad & A x \le b, \\
& A_{eq} x = b_{eq}, \\
& \ell \le x \le u
\end{aligned}\)
& Core relaxation form for simplex, IPM-LP, PDLP, HiGHS, and Gurobi. \\
\texttt{QP} & \texttt{QPModel} &
\(\begin{aligned}
\min/\max \quad & \tfrac12 x^\top Q x + c^\top x \\
\text{s.t.}\quad & A x \le b, \\
& A_{eq} x = b_{eq}, \\
& \ell \le x \le u
\end{aligned}\)
& Convex LCQP interface for native QP IPM and Gurobi. \\
\texttt{NLP} & \texttt{NLPModel} &
\(\begin{aligned}
\min/\max \quad & f(x) \\
\text{s.t.}\quad & g(x)=0, \\
& h(x) \le 0, \\
& \ell \le x \le u
\end{aligned}\)
& Generic nonlinear optimization contract with callback and optional symbolic forms. \\
\texttt{MILP} & \texttt{MIPModel} & LP plus
\(x_j \in \mathbb{Z}\) or \(x_j \in \{0,1\}\). & Native branch-and-cut input and external MILP adapter input. \\
\texttt{MINLP} & \texttt{MINLPModel} & NLP plus integer restrictions. & Native spatial B\&B shell and SCIP-facing model form. \\
\bottomrule
\caption{Canonical problem classes in the engine layer.}
\label{tab:engine-problem-classes}
\end{longtable}

The canonical model metadata are deliberately richer than just matrices and
vectors.
\begin{itemize}
  \item \texttt{VariableMeta} attaches type, lower bound, upper bound, and name to each decision variable.
  \item \texttt{SymExpr} and \texttt{SymbolicConstraint} allow an NLP to expose a symbolic representation when an external solver prefers exported algebra over callback-only evaluation.
  \item \texttt{MIPModel} carries integer index sets, an optional warm start, unit-commitment structure hints, and branching priorities, because the branch-and-cut engine uses all of them operationally.
\end{itemize}

\subsection{Validation and Dispatch Workflow}

The adapter path is intentionally thin.  Validation and numerical work are kept
separate so that front-end modules can fail early on malformed dimensions or
invalid bounds before any solver-specific data structure is allocated.

\paragraph{Algorithm E1: Canonical Adapter Dispatch}
\begin{algorithmic}[1]
\Procedure{\texttt{dispatch\_problem}}{$cls, model, registry, preferred$}
  \State $report \gets \texttt{validate}(model)$
  \If{$report.valid = \textbf{false}$}
    \State \Return unsuccessful result with validation errors
  \EndIf
  \If{$preferred \neq \emptyset$}
    \State $adapter \gets \texttt{registry.find\_by\_name}(preferred)$
  \Else
    \State $adapter \gets \texttt{registry.first\_for}(cls)$
  \EndIf
  \If{$adapter = \varnothing$}
    \State \Return unsuccessful result with ``no adapter available''
  \EndIf
  \State Dispatch to the type-specific entry point
  \Statex \quad \texttt{solve\_le}, \texttt{solve\_nle}, \texttt{solve\_lp}, \texttt{solve\_qp}, \texttt{solve\_nlp}, \texttt{solve\_milp}, or \texttt{solve\_minlp}
  \State \Return \texttt{SolveResult}
\EndProcedure
\end{algorithmic}

At the code level, the division of responsibility is straightforward.
\begin{itemize}
  \item \texttt{SolverAdapter} provides default ``unsupported problem class'' results.
  \item Each concrete adapter overrides only the methods it really supports.
  \item \texttt{AdapterRegistry} is just an ordered list.  Selection is therefore deterministic: the first supporting adapter wins unless the caller explicitly asks for one by name.
\end{itemize}

This simple registry design is important for application code such as market
simulation and time-series PF, which register Gurobi first, HiGHS second, and a
tuned native fallback last.

\subsection{Solver Inventory and Current Implementation Status}

Table~\ref{tab:engine-solver-inventory} separates the solver \emph{interfaces}
from the solver \emph{implementations}.  This is the part most likely to reduce
maintainer confusion, because the engine now contains both fully operational
solvers and placeholder interfaces meant for future completion.

\begin{longtable}{L{0.22\linewidth}L{0.15\linewidth}L{0.20\linewidth}L{0.33\linewidth}}
\toprule
Solver family & Main class or entry point & Problem class & Current status in source \\
\midrule
\endhead
PF Newton kernel & \texttt{NewtonSolver} & direct \texttt{SolverData} path & Fully implemented power-flow-specific Newton kernel with Jacobian-pattern caching, sparse factorization reuse, line search, and bus-type switching. \\
Generic linear solve & \texttt{NativeLinearAdapter} & \texttt{LE} & Fully implemented thin wrapper over the sparse linear solver abstraction. \\
Generic NLE solve & \texttt{NativeNewtonAdapter} & \texttt{NLE} & Fully implemented regularized Newton method with backtracking. \\
Generic NLP solve & \texttt{NativeNLPAdapter} & \texttt{NLP} & Fully implemented penalty-based line-search method using callback gradients and Hessians. \\
Native NLP IPM interface & \texttt{NativeIPMAdapter} & \texttt{NLP} & Fully implemented sparse primal-dual predictor-corrector IPM with reusable KKT factorizations, shared reduced-KKT Schur solves for small equality blocks, sparse secant or finite-difference Hessian recovery when \texttt{hess} is absent, startup curvature caching across repeated solves, and IPOPT fallback for unrecovered failures. \\
Convex QP / LCQP & \texttt{NativeLCQPAdapter} & \texttt{QP}, \texttt{LP} & Fully implemented Mehrotra-style primal-dual IPM for convex QP and LP special cases. \\
High-performance LP IPM & \texttt{NativeIPMLPAdapter} & \texttt{LP} & Fully implemented predictor-corrector IPM with scaling, correctors, cached node solves, and optional crossover support. \\
Dual simplex & \texttt{solve\_lp\_with\_basis} & \texttt{LP} & Fully implemented bounded-form simplex with standard-form conversion, warm starts, sparse basis reuse, and dual-simplex reoptimization. \\
PDLP / PDHG & \texttt{NativePDLPAdapter} & \texttt{LP} & Fully implemented first-order LP solver for large sparse problems. \\
MILP / MINLP tree search & \texttt{solve\_milp\_bc}, \texttt{solve\_minlp\_bc}, \texttt{NativeBranchAndCutAdapter} & \texttt{MILP}, \texttt{MINLP} & MILP branch-and-cut is fully implemented and mature.  The MINLP path now uses the native NLP relaxation, carries multiplier-informed branch scores into child nodes, and remains backed by SCIP/IPOPT-style external fallbacks where available. \\
HiGHS wrapper & \texttt{HighsAdapter} & \texttt{LP}, \texttt{MILP} & Fully implemented external file-based wrapper around the HiGHS executable. \\
IPOPT wrapper & \texttt{IpoptAdapter} & \texttt{NLP} & Available when IPOPT is built in.  Used for external NLP backend selection. \\
SCIP wrapper & \texttt{ScipAdapter} & \texttt{MINLP} & Available when SCIP support is built in. \\
Gurobi wrapper & \texttt{GurobiAdapter} & \texttt{LP}, \texttt{QP}, \texttt{MILP} & Native C-API wrapper for high-end LP/QP/MILP solves. \\
\bottomrule
\caption{Engine solver inventory and actual implementation status.}
\label{tab:engine-solver-inventory}
\end{longtable}

\paragraph{A practical consequence for readers of the code}
When the repository mentions a ``native IPM'' for NLP, that phrase must be read
carefully.  The engine contains:
\begin{itemize}
  \item a working native IPM for \textbf{LP} in \texttt{ipm\_lp\_solver.cpp},
  \item a working native IPM for \textbf{convex QP} in \texttt{lcqp\_solver.cpp},
  \item a working native IPM for \textbf{general NLP} in \texttt{ipm\_solver.cpp}, with reduced-KKT and missing-Hessian recovery paths,
  \item and a separate working \texttt{NativeNLPAdapter} that is not a barrier IPM, but a penalty-and-line-search Newton-type method.
\end{itemize}

\paragraph{Current performance status}
The most relevant benchmark result is now the structured sparse chain NLP in
\texttt{benchmark/nlp\_ipm\_benchmark.cpp}.  On the current implementation,
the exact-Hessian native path beats IPOPT on all benchmarked sizes

\[
\text{Native/IPOPT time ratio} = 0.419\ (n=64),\quad 0.518\ (n=128),\quad 0.773\ (n=256),
\]

while the missing-Hessian path is now also faster than IPOPT across the same
benchmark family.  After the sparse secant plus finite-difference Hessian
recovery pass, shared reduced-KKT solves, and startup curvature caching, the
native missing-Hessian solve achieves

\[
\text{Native/IPOPT time ratio} = 0.340\ (n=64),\quad 0.444\ (n=128),\quad 0.677\ (n=256).
\]

So the main nonlinear performance story has changed materially: the engine is
no longer merely ``equipped with a native NLP IPM,'' but already ahead of the
external IPOPT baseline on this structured sparse benchmark in both the
exact-Hessian and missing-Hessian regimes.  The remaining engineering work is
therefore about broadening that win across more difficult NLP families, not
about filling a missing core implementation.

\subsection{Key Algorithm Workflows}

\subsubsection{Power-Flow-Specific Newton Kernel}

\texttt{NewtonSolver} is the main reusable PF engine kernel rather than a
generic optimization adapter.  It operates on \texttt{powerflow::SolverData},
builds a Jacobian context from AC/DC bus states and converter modes, and then
reuses sparse-pattern analysis aggressively across repeated solves.

\paragraph{Algorithm E2: \texttt{NewtonSolver.solve}}
\begin{algorithmic}[1]
\Procedure{\texttt{NewtonSolver.solve}}{$data, opt, init$}
  \State Initialize voltage magnitudes, angles, DC voltages, and control states
  \State Detect AC slack bus and one DC slack bus per DC island
  \State Build Jacobian context from current bus types and converter modes
  \If{cached Jacobian pattern matches \texttt{data} structure and context layout}
    \State Reuse analyzed sparse pattern and linear solver object
  \Else
    \State Rebuild Jacobian sparsity pattern and analyze it once
  \EndIf
  \For{outer PV/PQ or converter-mode adjustment loop}
    \For{$iter = 1$ to \texttt{max\_iter}}
      \State Evaluate mismatch vector and Jacobian
      \If{residual norm is below tolerance}
        \State Assemble final PF result and return
      \EndIf
      \State Factorize Jacobian, solve for Newton step, and clip step components
      \State Run globalization / backtracking line search
      \State Apply PV/PQ switching and converter control-mode updates when enabled
    \EndFor
  \EndFor
  \State Return non-converged result with profiling counters
\EndProcedure
\end{algorithmic}

\subsubsection{Generic Nonlinear Equation Solver}

\paragraph{Algorithm E3: \texttt{NativeNewtonAdapter.solve\_nle}}
\begin{algorithmic}[1]
\Procedure{\texttt{solve\_nle}}{$F(x), J(x), x_0$}
  \State Validate dimensions and callbacks
  \State $x \gets x_0$
  \For{$iter = 1$ to \texttt{max\_iter}}
    \State Evaluate residual $f \gets F(x)$
    \If{$\lVert f \rVert_\infty \le \epsilon$}
      \State Return converged result
    \EndIf
    \State Form Jacobian $J(x)$
    \For{a small regularization ladder $\lambda$}
      \State Solve \((J + \lambda I) \Delta x = -f\)
      \State Try backtracking steps $x + \alpha \Delta x$
      \If{a trial decreases residual norm}
        \State Accept step and continue outer loop
      \EndIf
    \EndFor
    \If{no acceptable step exists}
      \State Return failure: line search failed
    \EndIf
  \EndFor
  \State Return failure: maximum iterations reached
\EndProcedure
\end{algorithmic}

\subsubsection{Generic NLP Solver Used Today}

\paragraph{Algorithm E4: \texttt{NativeNLPAdapter.solve\_nlp}}
\begin{algorithmic}[1]
\Procedure{\texttt{solve\_nlp}}{$f, \nabla f, \nabla^2 f, g, h, x_0$}
  \State Validate callbacks and project $x_0$ into variable bounds
  \For{$iter = 1$ to \texttt{max\_iter}}
    \State Evaluate objective, gradient, and optional Hessian
    \State Evaluate equality and inequality violations
    \State Build penalized merit model
    \Statex \quad $\phi(x)=f(x)+\tfrac{\rho}{2}\lVert g(x) \rVert_2^2+\tfrac{\rho}{2}\lVert \max(h(x),0) \rVert_2^2$
    \State Form penalized gradient and Hessian approximation
    \If{stationarity and feasibility are small enough}
      \State Return converged result
    \EndIf
    \State Regularize Hessian and solve for a step \(\Delta x\)
    \State Run backtracking on the penalized merit function
    \State Project trial point into variable bounds
    \If{no acceptable trial point exists}
      \State Return failure
    \EndIf
  \EndFor
  \State Return best iterate with status \texttt{Max iterations reached}
\EndProcedure
\end{algorithmic}

\subsubsection{Convex QP and LCQP Solver}

\paragraph{Algorithm E5: \texttt{NativeLCQPAdapter.solve\_qp}}
\begin{algorithmic}[1]
\Procedure{\texttt{solve\_lcqp}}{$Q, c, A_{eq}, b_{eq}, \ell, u$}
  \State Apply scaling to variables, equalities, costs, and bounds
  \State Initialize primal variables at the center of finite bounds
  \State Initialize lower/upper slacks and dual multipliers to positive values
  \State Build the KKT sparsity pattern once and analyze it symbolically
  \For{$iter = 1$ to \texttt{effective\_max\_iter}}
    \State Compute stationarity, equality, slack, and complementarity residuals
    \State Compute \(\mu\) from slack-dual products
    \If{primal feasibility, dual feasibility, and gap are below tolerance}
      \State Unscale solution and return
    \EndIf
    \State Update only the diagonal-dependent part of the KKT matrix
    \State Factorize KKT numerically
    \State Choose centering parameter \(\sigma\)
    \State Solve KKT system for primal, dual, slack, and multiplier directions
    \State Compute fraction-to-boundary step sizes
    \State Update primal and dual variables
  \EndFor
  \State Return last iterate with non-converged status
\EndProcedure
\end{algorithmic}

\subsubsection{High-Performance LP Interior-Point Method}

\paragraph{Algorithm E6: \texttt{NativeIPMLPAdapter.solve\_lp}}
\begin{algorithmic}[1]
\Procedure{\texttt{solve\_ipm\_lp}}{$LP$}
  \State Convert the LP to the engine's merged internal form with box bounds, equality rows, and slack variables
  \State Optionally presolve, scale, and cache repeated node-solve data structures
  \State Initialize primal variables, slacks, and duals strictly inside the feasible box
  \For{$iter = 1$ to \texttt{max\_iter}}
    \State Compute primal residual, dual residual, and complementarity gap
    \If{all three satisfy tolerance}
      \State Export primal solution, duals, and bound multipliers; return
    \EndIf
    \State Form diagonal barrier scaling terms and factorize normal equations
    \State Compute affine predictor direction
    \State Compute affine step lengths and predicted \(\mu_{aff}\)
    \State Choose centering parameter \(\sigma\) from the affine-gap ratio
    \If{the affine step is already nearly central}
      \State Skip expensive corrector and reuse affine direction
    \Else
      \State Form Mehrotra corrector right-hand side and solve corrector system
    \EndIf
    \State Apply fraction-to-boundary step sizes and update primal/dual variables
  \EndFor
  \State Return failure or best available iterate
\EndProcedure
\end{algorithmic}

This solver is the native LP workhorse for repeated bound-change problems.  The
important engineering feature is not only the Mehrotra iteration itself, but
also the \texttt{prepare\_for\_node\_solves()} and cached-node-solve path used by
branch-and-cut workers.

\subsubsection{Dual Simplex and Standard-Form Reoptimization}

\paragraph{Algorithm E7: \texttt{solve\_lp\_with\_basis}}
\begin{algorithmic}[1]
\Procedure{\texttt{solve\_dual\_simplex}}{$LP, basis\_hint$}
  \State Convert the original LP to bounded standard form
  \Statex \quad add slack, surplus, and artificial columns as needed
  \If{a valid parent basis hint exists}
    \State Rebuild basis state, reduced costs, and basic solution from the hint
    \If{dual feasible but primal infeasible}
      \State Reoptimize by dual simplex
    \ElsIf{warm start is numerically rejected and cold start is allowed}
      \State Fall back to cold-start Phase I / Phase II
    \EndIf
  \Else
    \State Build crash basis and run cold-start Phase I / Phase II simplex
  \EndIf
  \State Persist basis information for children
  \State Return primal solution, basis, reduced costs, and dual information
\EndProcedure
\end{algorithmic}

The bounded-form conversion is central to the engine architecture.  Child nodes
usually change only bounds, so the matrix structure remains fixed while the
basis and bound-shift terms are reused.

\subsubsection{First-Order PDLP Solver}

\paragraph{Algorithm E8: \texttt{NativePDLPAdapter.solve\_lp}}
\begin{algorithmic}[1]
\Procedure{\texttt{solve\_pdlp}}{$LP$}
  \State Convert the LP into a matrix-and-box form suitable for PDHG iterations
  \State Build both CSC and CSR sparse views of the matrix
  \State Apply Ruiz row/column equilibration
  \State Build Pock--Chambolle diagonal step sizes and optional IP-PMM regularization
  \State Initialize primal and dual iterates and their running averages
  \For{$iter = 1$ to \texttt{max\_iter}}
    \State Optionally extrapolate the primal iterate
    \State Update dual variables by a row-wise sparse pass and box projection
    \State Update primal variables by a column-wise sparse pass and bound projection
    \State Update ergodic averages
    \If{it is a check iteration}
      \State Evaluate primal violation, dual stationarity, and relative gap
      \State Update the best current or averaged iterate
      \State Trigger adaptive restart if progress stagnates or periodic restart fires
      \State Adapt step gain, primal weight, and PMM regularization
      \If{current or averaged iterate satisfies tolerances}
        \State Return best iterate
      \EndIf
    \EndIf
  \EndFor
  \State Unscale the best iterate and return
\EndProcedure
\end{algorithmic}

\subsubsection{Mixed-Integer Branch-and-Cut}

\paragraph{Algorithm E9: \texttt{solve\_milp\_bc}}
\begin{algorithmic}[1]
\Procedure{\texttt{solve\_milp\_bc}}{$MIP, opt$}
  \State Validate the model and run presolve / bound tightening
  \State Solve the root LP relaxation
  \State Run root cut rounds, primal heuristics, and incumbent discovery
  \If{the root solution is integral or the gap is closed}
    \State Return immediately
  \EndIf
  \If{parallel mode is enabled}
    \State Launch heterogeneous worker threads and shared queues
  \Else
    \State Start sequential best-first / depth-first / hybrid tree loop
  \EndIf
  \While{time, node, and gap limits are not hit and the node pool is not empty}
    \State Select a node and propagate local bounds
    \State Solve the node LP relaxation, preferably by simplex warm start
    \If{the node is infeasible or dominated by incumbent}
      \State Prune node
    \ElsIf{the node solution is integral}
      \State Update incumbent and global bound
    \Else
      \State Choose branching variable using pseudocost / reliability / fallback rule
      \State Create up and down children
      \State Optionally add tree cuts or conflict information
      \State Push surviving children back to the search structure
    \EndIf
  \EndWhile
  \State Return incumbent, best bound, gap, and detailed B\&C statistics
\EndProcedure
\end{algorithmic}

The full MILP machinery is larger than this summary.  The detailed cut families,
warm-start mechanics, presolve propagation, and parallel explorer design are
documented further in Section~\ref{sec:native-branch-and-cut}.

\subsubsection{External Backend Wrapper Pattern}

\paragraph{Algorithm E10: External Adapter Shell}
\begin{algorithmic}[1]
\Procedure{\texttt{solve\_external}}{$model, executable$}
  \State Validate the canonical model
  \State Export the model to an interchange format or backend API representation
  \State Launch the external solver or native C API call
  \State Parse objective, status, primal solution, and available dual data
  \State Map backend status to \texttt{SolveStats}
  \State Return \texttt{SolveResult}
\EndProcedure
\end{algorithmic}

In the current code base, HiGHS mainly follows the file-export pattern, while
Gurobi is linked through its C API.  IPOPT and SCIP are enabled only when their
build flags are available.

\subsection{Assessment Against the Self-Contained Computation-Engine Target}

Relative to the architecture goal in \texttt{ARCHITECTURE.md}, the current
engine is best described as a \emph{strong LP/MILP-centered computation engine
with reusable nonlinear interfaces}, not yet as a fully self-contained solver
stack across linear equation, nonlinear equation, MILP, and MINLP classes.

\subsubsection{Robustness Assessment}

\paragraph{Current strengths}
The framework is structurally sound in four important ways.
\begin{itemize}
  \item \textbf{Canonical contracts are clean.}  \texttt{problem\_types.hpp}, \texttt{problem\_validation.cpp}, and the adapter interface separate model definition, validation, dispatch, and numerical work in a maintainable way.
  \item \textbf{LP/MILP solve paths defend themselves well.}  The MILP path includes validation, presolve, warm starts, fallback logic, detailed statistics, and multiple root/node strategies instead of a single brittle algorithm.
  \item \textbf{PF and linear algebra kernels are reusable.}  \texttt{NewtonSolver}, sparse linear solves, and the LP/QP kernels have enough internal structure to serve as an actual computation-engine layer rather than one-off application code.
  \item \textbf{Compatibility wrappers reduce migration risk.}  The canonical home is now \texttt{hacdcpf::engine}, while older \texttt{core}/\texttt{solver} include paths still exist as shims for the rest of the code base.
\end{itemize}

\paragraph{Current risks and gaps}
Three issues materially limit the claim of a self-contained engine.
\begin{itemize}
  \item \textbf{Native MINLP is not currently complete.}  The MINLP B\&B driver calls \texttt{solve\_nlp\_relaxation()}, and that routine currently instantiates \texttt{NativeIPMAdapter}; however, \texttt{ipm\_solver.cpp} still returns ``not implemented'' for the native NLP-IPM solve.  In other words, the tree shell exists, but the native continuous relaxation underneath is incomplete.
  \item \textbf{The generic NLP path is usable but algorithmically modest.}  \texttt{NativeNLPAdapter} is a penalty-plus-line-search Newton method, not a full primal-dual barrier solver with robust filter/merit updates, restoration, or strong multiplier handling.  That is adequate for small smooth models, but it is a weaker foundation for large constrained NLP and for MINLP relaxations.
  \item \textbf{Regression protection is uneven on the hardest path.}  The type-system and adapter tests are healthy, but the branch-and-cut test target is not currently a clean safety net for the full engine because it can abort before reaching later scenarios.  That weakens confidence in day-to-day robustness of the mixed-integer core.
\end{itemize}

\subsubsection{Performance Assessment}

\paragraph{Where the engine is strong today}
The performance architecture is already credible in the LP/MILP center.
\begin{itemize}
  \item \textbf{Root LP performance is strong.}  The native LP-IPM path is heavily engineered with scaling, cached symbolic work, and node-solve preparation; internal benchmarks have shown it can outperform the native simplex root solve by a wide margin on large UC relaxations.
  \item \textbf{MILP tree engineering is real, not superficial.}  The branch-and-cut stack contains presolve, cut pools, pseudocost and reliability branching, warm-started simplex reoptimization, heuristics, and a heterogeneous parallel mode.
  \item \textbf{Sparse factorization reuse is treated as a first-class concern.}  This is visible across \texttt{newton\_solver.cpp}, \texttt{dual\_simplex.cpp}, \texttt{ipm\_lp\_solver.cpp}, and the sparse LU support files.
\end{itemize}

\paragraph{Main bottlenecks versus the target architecture}
The remaining bottlenecks are also clear.
\begin{itemize}
  \item \textbf{Large-case MILP competitiveness is still uneven.}  Internal benchmarking shows the native B\&C path can be competitive on some UC families, but on harder large cases the external commercial or dedicated MILP solvers still retain a sizable advantage, mainly from stronger tree closure and LP reoptimization throughput.
  \item \textbf{External-wrapper overhead remains in practical workflows.}  HiGHS and SCIP mainly use file export plus subprocess execution, which is fine for portability but weaker than an in-memory backend when low latency or high-frequency repeated solves are required.
  \item \textbf{Nonlinear performance headroom is large.}  Because the native NLP-IPM path is missing, the engine does not yet have a high-performance native continuous relaxation for general MINLP.  This is the main architectural hole between today's implementation and the intended computation-engine role.
\end{itemize}

\subsubsection{Readiness Summary}

For the four solver families named in the architecture document, the present
readiness is:
\begin{center}
\begin{tabular}{L{0.18\linewidth}L{0.22\linewidth}L{0.46\linewidth}}
\toprule
Family & Current readiness & Comment \\
\midrule
Linear equation & Strong & Ready as a reusable native engine kernel \\
Nonlinear equation & Strong & Newton-style solve path is available and robust for structured models \\
MILP & Strong--Moderate & Native path is mature, but still not uniformly best-in-class on large hard instances \\
MINLP & Partial & Interface and tree shell exist, but native continuous relaxation is not yet complete \\
\bottomrule
\end{tabular}
\end{center}

Therefore, the most accurate near-term roadmap is not to redesign the engine
abstraction, but to complete the missing native NLP/MINLP core under the
existing abstraction: a real NLP-IPM or SQP-quality relaxation solver, then
MINLP-specific reliability work on branching, node recovery, and regression
testing.

\subsection{Recommended Reading Order for Maintainers}

For someone trying to understand or extend the engine, the lowest-friction order
is:
\begin{enumerate}
  \item \texttt{problem\_types.hpp} and \texttt{problem\_validation.cpp},
  \item \texttt{solver\_adapter.hpp} and \texttt{adapter\_registry.cpp},
  \item \texttt{native\_adapters.cpp} to see which kernels are actually selected,
  \item one kernel family in depth: \texttt{dual\_simplex.cpp}, \texttt{ipm\_lp\_solver.cpp}, \texttt{lcqp\_solver.cpp}, or \texttt{pdlp\_solver.cpp},
  \item \texttt{branch\_and\_cut.cpp} and its helper files for mixed-integer logic,
  \item \texttt{external\_adapters.cpp} only after the canonical forms and native solve path are already clear.
\end{enumerate}

That order mirrors the architecture itself: understand the problem contracts
first, then the dispatch layer, then the heavy numerical kernels.